\scanpage{27}
1.3. Notons que $\unsure{\mathfrak S_3|\mathbb G_3}$ opère sur $C_n$ de façon évidente,
et les $C_{n'}\longrightarrow C_n$ sont des $\unsure{\mathfrak S_3|\mathbb G_3}$-morphismes,
donc $\unsure{\mathfrak S_3|\mathbb G_3}$ opère sur $C_\infty$.
L'opération de $\unsure{\mathfrak S_3|\mathbb G_3}$ sur $C_1$ \inserted{$\simeq\mathbb P^1_{\mathbb Q}$}
est celle déduite des opérations de $\unsure{\mathfrak S_3|\mathbb G_3}$ sur $\{0,1,\infty\}$.
\inserted{Enfin, $\unsure{\mathfrak S_3|\mathbb G_3}$ opère sur $\Pi_n$, de façon que
$\unsure{\mathfrak S_3|\mathbb G_3}$ opère sur $(C_n,\Pi_n\ldots)$ -- donc
$\unsure{\mathfrak S_n|\mathfrak S_3}$ opère sur $(C_\infty,\Pi_\infty)\ldots$}

1.4. Considérons la clôture alg. $\overline{\mathbb Q}$ de $\mathbb Q$ dans $\mathbb C$,
d'où par changement de base des schémas $\overline C_n$, $\overline C_n^*$,
$\overline C_\infty$, $\overline C_\infty^*$ sur $\overline C_1=\mathbb P^1_{\overline{\mathbb Q}}$
et $\overline C_1^*=\mathbb P^1_{\overline{\mathbb Q}}\setminus\{0,1,\infty\}$.
Sur ces schémas $C_n$
\deleted{$\unsure{\Gamma=\operatorname{Gal}}$}
\[
(7)\qquad\Gamma=\operatorname{Gal}(\overline{\mathbb Q}/\mathbb Q)
\]
opère de façon évidente, et \unsure{ajoutant|adjoignant} les groupes
\[
(8)\qquad\left\{\begin{aligned}
&\inserted{\pi_n=}\deleted{\unsure{\Gamma}}\,
\Pi_n(\overline{\mathbb Q})=
\unsure{\overline\mu_n^{\,3}|\mu_n^3}/\unsure{\overline\mu_n|\mu_n}
\simeq\unsure{\overline\mu_n|\mu_n}\otimes_{\mathbb Z}\mathbb Z^3/\mathbb Z\\
&\hspace{4cm}\simeq\unsure{\overline\mu_n|\mu_n}\otimes_{\widehat{\mathbb Z}}
(\widehat{\mathbb Z}^{\,3}/\widehat{\mathbb Z})\\
&\text{où}\quad\inserted{\mu_n=\mu_n(\overline{\mathbb Q})}\\
&\qquad\deleted{\unsure{\overleftarrow\mu_\infty|\Pi_\infty}(\overline{\mathbb Q})=
\varprojlim_n\mu_n(\overline{\mathbb Q})}
\end{aligned}\right.
\]
et $\unsure{\mathfrak S_3|\mathbb G_3}$. On voit tout de suite que ces opérations
s'ajustent de façon à ce qu'opère sur ces courbes le groupe
\[
(\unsure{9|8})\qquad\mathcal E_n=
(\unsure{\mathfrak S_3|\mathbb G_3}\times\Gamma)
\mathbin{\underset{\scriptscriptstyle 1/2}{.}}\pi_n
\]
% archive foliation: 26
